\section{Firm Theory and the Agency Problem}\label{sec:firms}

The framework developed in the preceding sections treats firms as passive executors of
the planner's output vector $\vec{X}_t^*$. This section relaxes that assumption by
introducing \textit{active firms} that choose their own production targets, report them
to the planner, and face an incentive mechanism that aligns their private optimum with
the socially optimal plan. The resulting structure provides a decentralised
implementation of the central plan that does not require any firm to observe the global
input-output structure.

\subsection{Firm Income and the Endogenous Tax Rate}\label{sec:firm-income}

Each sector $j = 1, \ldots, n$ receives income proportional to its value added net of
the tax rate $\tau_t$:

\begin{equation}\label{eq:firm-income}
    \Pi_{f,j,t} = (1 - \tau_t)\, v_{j,t}\, X_{j,t},
\end{equation}
where $v_j$ is the value-added coefficient and $X_{j,t}$ is actual gross output. The tax
rate $\tau_t$ funds both government expenditure $\vec{G}_t$ and net capital formation
$\Delta\vec{K}_t$. From the shadow-price-weighted fiscal identity:

\begin{equation}\label{eq:fiscal-identity}
    \tau_t\;\vec{\pi}_t \cdot \vec{F}_t
        = \vec{\pi}_t \cdot \bigl(\Delta\vec{K}_t + \vec{G}_t\bigr),
\end{equation}
where $\vec{F}_t = \vec{C}_t + \Delta\vec{K}_t + \vec{G}_t$ is aggregate final demand
and $\vec{\pi}_t \cdot \vec{F}_t = \mu_{t }\vec{v}_{t}^\top \vec{X}_t^*$ equals nominal GDP. It
follows that $\tau_t$ is \textit{endogenously determined} each quarter residually:

\begin{equation}\label{eq:tau-endogenous}
    \boxed{\tau_t = \frac{\vec{\pi}_t \cdot \bigl(\Delta\vec{K}_t + \vec{G}_t\bigr)}
                  {\vec{\pi}_t \cdot \vec{F}_t}.}
\end{equation}

\begin{remark}[Policy Degrees of Freedom]\label{rem:policy}
Equation~\eqref{eq:tau-endogenous} shows that $\tau_t$ is not a free policy
parameter. Once $\vec{X}_t^*$ is determined by the optimization problem and the
political process has set $\vec{G}_t$, the tax rate adjusts automatically to fund
implied investment and government expenditure. The \textit{only} genuinely exogenous
policy inputs to the system are the government expenditure vector $\vec{G}_t$ and the
subsistence floor $\vec{\gamma}$. Every other macroeconomic variable is computed.
\end{remark}

\begin{assumption}[Feasibility of Tax Rate]\label{ass:tau-feasible}
We impose $0 < \tau_{\min} \leq \tau_t < 1$ for all $t$, where $\tau_{\min} > 0$ is a
small floor ensuring the government can always fund $\vec{G}_t$. The upper bound is
guaranteed by the material balance whenever households consume a strictly positive bundle.
\end{assumption}

\subsection{The Self-Reporting Mechanism}\label{sec:self-report}

Rather than receiving an administratively imposed target, each firm $j$ observes the
planner's announced gross output vector $\vec{X}_t^*$ and the unit income
$(1-\tau_t)v_j$, then independently \textit{reports} an intended production quota
$X_{j,t}^*$. This is a two-stage process:

\begin{enumerate}
    \item \textbf{Report stage:} The firm chooses and submits $X_{j,t}^*$ --- its
          intended output for the period.
    \item \textbf{Verification stage:} Actual output $X_{j,t}^{\text{actual}}$ is
          recorded bilaterally from what counterparties receive, not from the firm's
          own claim. This makes the performance measure tamper-evident.
\end{enumerate}

The separation of the reported quota from the verified actual output is the
institutional feature that gives the incentive mechanism its force.

\subsection{The Incentive Function}\label{sec:incentive-function}

Total firm income combines the linear value-added payment with a concave
bonus-penalty term:

\begin{equation}\label{eq:total-firm-income}
    \Pi_{f,j,t}
        =  \left[ 1 + \phi_j \cdot f(r) \right] (1-\tau_t)v_j X_{j,t}^{\text{actual}}
\end{equation}
where $\phi_j > 0$ is a sector-specific bonus scale and $f : \mathbb{R}_{++} \to
\mathbb{R}$ satisfies the following:

\begin{assumption}[Incentive Function Properties]\label{ass:incentive}
\hfill
\begin{enumerate}
    \item $f(1) = 0$ \quad --- no bonus or penalty at exact quota;
    \item $f'(r) > 0$ for all $r > 0$ \quad --- more output always raises income;
    \item $f''(r) < 0$ for all $r > 0$ \quad --- strictly concave; diminishing
          marginal reward for overproduction;
    \item $f(r) \to -\infty$ as $r \to 0^+$ \quad --- catastrophic penalty for
          near-zero output;
    \item $\sup_{r>1} f(r) < \infty$ \quad --- overproduction bonus saturates.
\end{enumerate}
\end{assumption}

A canonical specification satisfying Assumption~\ref{ass:incentive} is:

\begin{equation}\label{eq:incentive-log}
    f(r) = \ln r - (r - 1),
\end{equation}
which has a unique global maximum of zero at $r = 1$, is strictly negative everywhere
else, and diverges to $-\infty$ as $r \to 0^+$.

\begin{remark}[Asymmetric Loss Structure]\label{rem:asymmetric}
The strict concavity of $f$ implies that underproducing by $\varepsilon$ relative to
quota incurs a strictly larger income loss than the income gain from overproducing by
the same $\varepsilon$. This asymmetry pushes rational firms toward or slightly above
their announced quota.
\end{remark}

\subsection{Incentive Compatibility: Truthful Quota Reporting}\label{sec:incentive-compatibility}

\begin{theorem}[Truthful Reporting]\label{thm:truthful}
For any anticipated actual output $X_{j,t}^{\text{actual}} = x > 0$, the unique
income-maximising reported quota is $X_{j,t}^* = x$. Firms truthfully report their
intended production.
\end{theorem}

\begin{proof}
Fix $x > 0$. Income as a function of the reported quota $s > 0$ is:
\begin{equation*}
    \Pi(s) = (1  + \phi_j\, f(x/s) ) (1-\tau_t)\,v_j\,x .
\end{equation*}
Since $\phi_j > 0$, maximising $\Pi$ over $s$ is equivalent to maximising $f(x/s)$.
Setting $r = x/s$, this reduces to maximising $f(r)$ over $r > 0$. By
Assumption~\ref{ass:incentive}(iii) and $f'(1) = 0$, the unique maximiser is $r^* = 1$,
i.e.\ $s^* = x$.
\end{proof}

The mechanism is therefore incentive-compatible by design: misreporting the quota
(high or low) strictly reduces the firm's own income. No external enforcement of
honesty is required. Though the system has a single free parameter $\phi$ which has to be determined administratively.

\subsection{Resolution of the Soft Budget Constraint}\label{sec:soft-budget}

Kornai~(1992) identified the \textit{soft budget constraint} as the central pathology
of classical planned economies: firms could miss targets and expect state bailouts,
rendering any penalty scheme non-credible. Three structural features of the present
mechanism eliminate this pathology without requiring any discretionary intervention by
the planner.

\begin{enumerate}
    \item \textbf{Self-chosen quotas.} Because firms set their own $X_{j,t}^*$,
          they cannot appeal to the planner on the grounds that targets were
          unrealistic or administratively imposed. The quota is the firm's own
          commitment device against itself.
    \item \textbf{Bilateral output verification.} Actual output is measured from
          what counterparties receive. The firm cannot manipulate the performance
          measure ex post.
    \item \textbf{Automatic penalty application.} The bonus-penalty
          formula~\eqref{eq:total-firm-income} is executed algorithmically. No
          human decision to apply or waive a penalty is required; the mechanism
          is self-enforcing.
\end{enumerate}

\begin{proposition}[Hard Budget Constraint]\label{prop:hard-budget}
Under Assumption~\ref{ass:incentive} and bilateral output verification, no Nash
equilibrium of the firm reporting game involves systematic underproduction relative
to the reported quota. The budget constraint is hard by construction.
\end{proposition}

\begin{proof}
Suppose firm $j$ anticipates $X_{j,t}^{\text{actual}} < X_{j,t}^*$. By
Theorem~\ref{thm:truthful}, the firm should have reported $X_{j,t}^* =
X_{j,t}^{\text{actual}}$, obtaining $f(1) = 0$ rather than $f(r) < 0$ for $r < 1$.
Any strategy resulting in underproduction relative to the reported quota yields
strictly lower income than truth-telling. No firm prefers systematic underproduction
in equilibrium.
\end{proof}

\subsection{Decentralised Implementation of the Social Optimum}\label{sec:decentralised}

The firm-level mechanism constitutes a \textit{decentralised implementation} of the
central plan in the sense of Kantorovich~(1965): each firm independently chooses the
socially optimal output without observing the global IO structure. This provides a
stronger answer to Hayek's~(1945) knowledge problem than shadow prices alone. The
knowledge problem is addressed not only by communicating scarcity through prices, but
by allowing firms to \textit{reveal capacity constraints} through their reports ---
information that is genuinely private and unavailable to the planner's IO model.

\begin{corollary}[Lange--Kantorovich Decentralisation]\label{cor:lange}
At the equilibrium of the firm reporting game, the aggregate of individually optimal
firm outputs equals the socially optimal output vector $\vec{X}_t^*$. The
decentralised equilibrium implements the social optimum.
\end{corollary}

\begin{proof}
The primal problem~\eqref{eq:planning_problem} is strictly convex: the objective
$\hat{W}_t(\vec{C})$ is strictly concave and the feasible set is a convex
polytope. Since the interior of the feasible set is non-empty (Slater's
condition holds whenever $\vec{K}_t \gg B\vec{X}_t$ and $L > \vec{l}^\top
\vec{X}_t$), strong duality obtains and there is no duality gap
\citep{boyd2004convex}.

At the optimal dual variables $(\vec{\lambda}_K^*, \lambda_L^*)$, the
Lagrangian~\eqref{eq:lagrangian} decomposes into independent firm-level
subproblems. Differentiating with respect to $\vec{X}_t$ and applying the
first-order condition~\eqref{eq:foc_production} yields
\begin{equation}
    (I - \bar{A})^\top \vec{\pi} \;=\; \lambda_L \vec{l} \;+\; B^\top
    \vec{\lambda}_K,
\end{equation}
so that each firm $f$'s income-maximisation problem
\begin{equation}
    \max_{\vec{x}_f \geq \vec{0}} \;\vec{v}_\text{eff}^\top \vec{x}_f
    \quad \text{s.t.} \quad B\vec{x}_f \leq \vec{K}_f,
\end{equation}
is precisely the Lagrangian subproblem obtained by decomposing~\eqref{eq:planning_problem}
at the optimal dual variables, where $\vec{v} = (I - \bar{A})^\top\vec{\pi}$
is the value-added coefficient vector. To see this, note that
\begin{multline}
    \vec{\pi} \cdot \vec{X}_f - \vec{\pi} \cdot (\bar{A}\vec{X}_f) \\
    \;=\;  \bigl((I-\bar{A})^\top\vec{\pi}\bigr) \cdot \vec{X}_f\\
    \;=\; \bigl(\lambda_L \vec{l} + B^\top\vec{\lambda}_K\bigr)
          \cdot \vec{X}_f\\
    \;\propto\; \vec{v}_\text{eff} \cdot \vec{X}_f,
\end{multline}

where the second equality substitutes the FOC and holds at the optimum of
the dual ascent procedure. By the strong duality theorem, the aggregate of
the individual firm solutions recovers the primal optimum:
\begin{equation}
    \sum_f \vec{x}_f^{\,*} \;=\; \vec{X}_t^*.
\end{equation}
The decentralised equilibrium therefore implements the social optimum.
Unless and until there is not degeneracy of the optimum points for the firms collectively.
\end{proof}